\section{\raggedright Statement language and reproducibility}
The Companion uses QAM as an object and assertion language, MML for
typed instructions and MMA for their execution semantics. For the
physical arguments in this paper, their essential role is to keep an
assertion attached to its domain, assumptions, calculation and remaining
interpretation requirements. A guard denotes a condition still needed
for a stated inference. A finite calculation and a general proof have
different scopes, even when both concern the same formula.

The projection identity provides a complete elementary example. Its
inputs are a Hilbert space, $J^2=I$, $J^*=J$ and the definition of $C$;
the calculation in the preceding section proves $C^2=C$. An exact
finite instance is
\begin{equation}
 P=\frac12\begin{pmatrix}1&1\\1&1\end{pmatrix},\quad
 M=\frac14\begin{pmatrix}5&3\\3&5\end{pmatrix},\quad
 N=\begin{pmatrix}1&0\\0&2\end{pmatrix}.
\end{equation}
Direct multiplication gives $P^2=P$, $[M,P]=0$, and
\begin{equation}
 [N,P]=\frac12\begin{pmatrix}0&-1\\1&0\end{pmatrix}\ne0.
\end{equation}
An exact rational check therefore confirms the projection and one
commuting transport while rejecting the other transport as a commuting
example. This rejection is a control result, not a failure of the
projection theorem. None of these matrix equations identifies a native
field, a physical transport or a mass operator.

Ordinary mathematical notation is retained throughout the main text.
The accompanying assertion records specify precise source versions,
prerequisites and finite checks; they allow the chain of reasoning to
be inspected without treating content identity as proof. The original
mass-readout arguments are collected in Appendix~A, while the separate
audit records retain the detailed provenance and unresolved obligations.
\subsection{An assertion is more than a displayed formula}
Consider the familiar expression $C^2=C$. The formula is ambiguous
until the domain of $C$, the scalar field, the composition operation
and the assumptions on its defining involution are fixed. It also
does not state whether the assertion is adopted, derived generally,
or checked only for a particular matrix. A physical application adds
the question of what is measured by the range of $C$.

For exposition, write an assertion record as
\begin{equation}
 r=(\varphi,\mathcal F,s,W,G),
\end{equation}
where $\varphi$ is the statement, $\mathcal F$ its declared context,
$s$ its evidence status, $W$ its sources and premises, and $G$ its
remaining guards. This tuple summarizes the record's role here; it
is not an additional literal MML instruction. Two records with the
same displayed formula can describe different assertions when their
contexts or premises differ.

QAM is the project's language for such objects and assertions. MML
specifies operations on the declared objects. MMA interprets the
supported instructions and returns execution information. The
implementation used in the examples is a finite interface, not an
unrestricted interpreter for every symbolic statement in the Companion.
Consequently an indexed formula is not automatically executable, and
a successful execution is not automatically a proof of its physical
interpretation.

An object can carry explicit index-role data. The implemented interface
records an index label, its side, variance, role and declared space.
The side and role must be a permitted pair in that interface. A
contraction must also specify which spaces are paired. For instance,
two upper structural indices do not define a scalar product until
a compatible comparison tensor is supplied. Merely printing repeated
labels cannot replace that tensor. This is especially relevant when
an internal comparison is proposed as a spacetime contraction.

\subsection{Derivations, combinations and frame changes}
Three simple operations illustrate the distinction between manipulating
records and proving their contents. A derivation step attaches a
conclusion to its premise record. A join combines records in a common
context. A conversion transports a record through an explicitly
declared map. In the independent provenance checker these appear as
the constructors \texttt{step}, \texttt{join} and \texttt{conv}.

The checker uses an ordered status vocabulary: open, reading,
conditional, computed and proved. This is a protocol ordering used
to prohibit unsupported promotion. It is not a universal ranking of
scientific confidence. In particular, a conditional general theorem
and an exact computation of one example answer different questions.
A valid record tree cannot obtain a stronger status simply by
combining or converting its premises. Its witness inputs remain
tracked, and a frame conversion must name a declared comparison map.

These structural properties have formal record-level statements in
the project's Lean implementation. Their scope is important. A
theorem that well-formed records do not lose witness identifiers
does not prove that a witness describes the correct physical system.
Similarly, a declared map can be type-correct before its metric or
observable-preservation properties have been proved. The semantic
conditions must appear as their own premises or guards.

In the LaTeX interface, a frame declaration introduces an identifier,
a map declaration specifies its source and target, and a formula
declaration supplies a formula, context and status. The actual macro
names include \texttt{qamframe}, \texttt{qammap},
\texttt{qamformula} and \texttt{qamphysconvert}. These are document
commands rather than assertions that a physical conversion exists.
They allow the same ordinary equation to be displayed with its
explicit profile when taken outside its local context. A laboratory
formula and a formula in a boosted-observer profile must therefore
retain the information identifying their different preparation and
comparison data.

\subsection{A complete physical example: information lost by a readout}
Take classical electromagnetic fields in units $c=1$ and form the
three-by-two matrix $F=[\boldsymbol E\mid\boldsymbol B]$. On nonzero
fields define the normalized Gram readout
\begin{equation}
 \rho(F)=\frac{F^{\mathsf T}F}{\operatorname{Tr}(F^{\mathsf T}F)}.
\end{equation}
This positive two-by-two matrix records relative norms and the scalar
product of the two vectors. It discards their common scale and their
orientation in physical space. These omissions matter under an
observer transformation whose direction is specified in the laboratory.

For a boost of velocity $v$ in the positive $z$ direction, consider
the pair of preparations
\begin{equation}
 \boldsymbol E=2e_x,\qquad \boldsymbol B=b e_y,\quad b=\pm1.
\end{equation}
Both have the same input register $\rho=\operatorname{diag}(4/5,1/5)$.
The ordinary field transformation gives
\begin{equation}
 \boldsymbol E'_b=\gamma(2-vb)e_x,\qquad
 \boldsymbol B'_b=\gamma(b-2v)e_y,
 \quad \gamma=(1-v^2)^{-1/2}.
\end{equation}
At $v=1/4$, the common factor cancels from the normalized Gram
readout and the two outputs are
\begin{equation}
 \rho(F'_{+1})=\operatorname{diag}(49,4)/53,\qquad
 \rho(F'_{-1})=\operatorname{diag}(9,4)/13.
\end{equation}
Thus there is no single-valued boosted readout depending only on
the original normalized register for this family. A map on that
register would need to assign two different outputs to one input.
This is an elementary failure of information sufficiency, not a
failure of the Lorentz transformation on the full fields.

The example explains why a frame label alone cannot repair a reduced
description. One must either retain additional preparation data or
replace a single predicted value by the set of outputs compatible
with the data actually available. This same distinction will later
appear in the hidden-lift term of projected dynamics.

\subsection{An exact interval prediction from an additional measurement}
Allow the laboratory preparation to vary as
\begin{equation}
 \boldsymbol E=2e_x,\qquad
 \boldsymbol B=t e_y+u e_z,\qquad t^2+u^2=1.
\end{equation}
All these fields again have the same input Gram register. Suppose
an extra measurement constrains the signed laboratory component
$t=\boldsymbol B\cdot e_y$ to $[a,b]\subseteq[-1,1]$. At the
same boost speed $v=1/4$,
\begin{equation}
 \boldsymbol E'=\gamma(2-t/4)e_x,\qquad
 \boldsymbol B'=\gamma(t-1/2)e_y+u e_z.
\end{equation}
The vectors remain orthogonal. Divide their squared norms by
$\gamma^2$ and set $A(t)=(2-t/4)^2$. Then
\begin{equation}
 \frac{|\boldsymbol B'|^2}{\gamma^2}
 =(t-1/2)^2+(15/16)(1-t^2)=A(t)-45/16.
\end{equation}
Consequently the first diagonal entry is
\begin{equation}
 p(t)=\frac{A(t)}{2A(t)-45/16}
     =\frac{(8-t)^2}{2(8-t)^2-45}.
\end{equation}
The denominator is strictly positive on the admitted interval.
Moreover $A$ decreases with $t$ and
$dp/dA=-(45/16)/(2A-45/16)^2<0$, so $p$ increases. The exact
image of the measurement interval is therefore
\begin{equation}
 \rho'\in\{\operatorname{diag}(p,1-p):p(a)\leq p\leq p(b)\}.
\end{equation}
Every intermediate value is attained by some preparation in the
specified family. This is a deterministic interval statement; no
probability distribution on $t$ has been assumed.

For $[a,b]=[-1/4,1/4]$ the exact result is
\begin{equation}
 p\in[121/162,\ 961/1202]
      \subset[9/13,\ 49/53].
\end{equation}
The wider interval is the output when only $t\in[-1,1]$ is known.
The additional measurement has therefore narrowed the prediction
without reconstructing an arbitrary orientation. The numbers in
this example are a controlled classical preparation, not fitted
lepton masses or an empirical confirmation of the common-carrier
hypothesis.

\subsection{From the assertion to the MML request and the MMA receipt}
The QAM record for this example declares a map
$\mathrm{Lab}\to\mathrm{Boosted}$ and two conditional assertions:
the preparation, and the interval rule derived from it. In the
implemented request the instruction is \texttt{READOUT}. Its
arguments name the recipe \texttt{em.gram-boost-range}, the preparation
\texttt{E2x-B-unit-yz}, source frame \texttt{Lab}, target frame
\texttt{Boosted}, boost \texttt{z:1/4}, and endpoints
\texttt{-1/4} and \texttt{1/4}. The endpoints are rational strings,
so the adapter can evaluate them without a floating rounding step.

The stored receipt for request \texttt{gram\_boost\_interval} returns
the exact target interval $[121/162,961/1202]$. Its recorded premise
chain links the interval rule to the stated laboratory preparation.
This receipt was retrieved by its content hash and checked through
the read-only evidence interface for this edition. The complete
identifier and source fingerprints belong in the accompanying evidence
files, where a reader can distinguish this saved calculation from
a newly executed or replayed one.

Its execution status is \texttt{continue}, with no outstanding
adapter obligations and no missing recorded premises. The premises
themselves remain conditional. In particular, the same receipt records
that closure is not proved and that no general mathematical relation
is certified by the interpreter. The analytic derivation above supplies
the general interval argument; the receipt verifies the specified
instance and its recorded execution context.

The other execution outcomes have equally limited meanings.
\texttt{guard} identifies an unmet prerequisite for the requested
operation, whereas \texttt{reject} identifies a request that cannot
be accepted under the interpreter's contract. These outcomes are not
statistical likelihoods or statements about the truth of an entire
physical theory. An empty list of additional request guards does not
erase the assumptions in the preparation or the rule.

\subsection{What is preserved when the document is regenerated}
The document store retains immutable assertion versions, source
identities and ordered selections. A reader edition can move an
original proof to an appendix while keeping its exact source
association. A change of prose order does not create a new inference,
and a content hash does not settle whether an imported theorem heading
is justified. The expanded physical argument must continue to state
its mathematical assumptions in the body, rather than relying on a
reader to discover them in an audit log.

This division of work is useful in a multi-branch programme. The main
paper explains a restricted set of connected results. Its evidence
package supplies the reproducible records for those results. The long
Companion retains alternative routes and their history, including
failed or qualified claims. A future change to a preparation or a
comparison tensor must be traced through dependent results before the
corresponding physical conclusion can be reused.
